\uuid{OzBi}
\titre{Équation droite/plan}
\auteur{bodin ; hanani ; mzari}
\organisation{exo7}
\licence{CC BY-NC-SA 4.0}
\datecreate{2020-01-20}
\niveau{L1}
\module{Géométrie}
\chapitre{Géométrie de l'espace}
\sousChapitre{Plans}
\difficulte{3}
\theme{Plans -- Droites, équation droite, plan}
\exo7qcmid{lille-1.geometrie-espace.017}

\contenu{
\texte{
Pour ces questions, l'espace est muni d'un repère orthonormé direct  $(O, \vec{i}, \vec{j}, \vec{k})$.
}

\begin{question}
\qclassification{141.03,141.04}
\qtags{motcle=équation droite/plan}
Soit $D$ la droite d'équations :  $ \left\{\begin{array}{ccl}x-y&=&-1\\y-z&=&0 \end{array}\right.$ et $P$ le plan contenant $D$ et parallèle à la droite d'équations  $D' : \left\{\begin{array}{ccl}x+z&=&0\\x-y&=&2 \end{array}\right.$.
 Une équation cartésienne de $P$ est :
\begin{answers}

       \bad{$x-z=1$}

      \bad{$x-y=0$}

     \bad{$y-z=0$}

     \good{$x-y=-1$}

\end{answers}
\begin{explanations}
$\vec{u}(1,1,-1)$ est un vecteur directeur de la droite $D'$ qui  n'appartient pas au  plan  $y-z=0$ .   Donc $P$ est différent du plan d'équation :  $y-z=0$ et donc une équation cartésienne de $P$ est  de la forme : $(x-y+1) + \alpha(y-z)=0, \, \alpha \in \Rr$. On calcule $\alpha$ de sorte que $\vec{u}(1,1,-1)$ soit un vecteur de  $P$.
\end{explanations}
\end{question}
}
